\subsection{Component propagation at the extra vertices in the rank-four construction}
\label{sec:count-three-components}

Having the required residual links describes a path of points. At its extra vertices one still needs to pass between components without losing the condition that both ordinary ideals vanish.

\begin{proposition}[Propagation through the extra vertices]\label{prop:count-three-extra-vertices}
At the extra vertices of the residual path in
Proposition~\ref{prop:count-three-residual-path}, use the
depth and inertia-coincidence dimension bounds of
Proposition~\ref{prop:tower-component-propagation}. If \(P,P'\) are any two component primes through such a vertex and \(J\) denotes the ordinary ideal there, 
\[
 J\subseteq P\quad\Longrightarrow\quad J\subseteq P'.
\]
 Along the
four links this gives the backward propagation from the integral
seed of Proposition~\ref{prop:count-three-ordinary-attachment}.
\end{proposition}

\begin{proof}
For the first three vertices of
Section~\ref{sec:count-three-chain}, the ordinary
component argument of Paragraph~\ref{par:count-two-propagation}
already applies. We prove the additional statement
for a residual datum in \(\mathcal A\) having global root
\(-w_\Psi\), or both \(-w_\beta,-w_\vartheta\).
The ideal in the argument is the same sum of the two
ordinary Hecke ideals.

\proofstep{The reducible loci at an extra vertex.}
Use the nine supports in the component argument.
The ninth is
\[
                   Y=\{\pm w_\beta,\pm w_\vartheta\}.
\]
For every other support meeting the extra vertex, the
tangent dimension is at most \(d+20C\).
Indeed, use the symplectic cocharacter with exponents
\((0,0,1,1)\). It sends the global matrices to the
diagonal and the frame at \(v\) to
\[
                     N_v=1+b_v(E_{10}-E_{32}),
                     \qquad b_v\ne0.
\]
The tangent dimension can only increase at this
limit. The first-column flag equation at a datum in
\(\mathcal A\) gives \(d_0=d_1\) on each marked
decomposition group, since global row one has no
off-diagonal entries and \((N_v)_{10}\ne0\).
The same equality holds at the limit.

For a global tangent cocycle \(U\) and a frame
variation \(V_v\), put
\(Z(s)=U(s)+\mathcal D(s)V_v\mathcal D(s)^{-1}-V_v\)
at the local group. Its conjugate by \(N_v\) is
upper triangular. Direct multiplication gives
\[
 Z_{10}=b_v^2 Z_{01}+b_v(Z_{00}-Z_{11}),\qquad
 Z_{21}=0,\qquad Z_{30}=0 .
\]
For the last equality use also entries \((2,0),(3,1)\)
of the upper triangular condition.
First fix the two independent global diagonal
cocycles. This costs at most \(2C\).
After fixing the two diagonal cocycles, the local equations give
\[
       U_{10}=b_v^2U_{01},\qquad [U_{21}]=[U_{30}]=0
         \text{ in the local root cohomology}.
\]
The root character for \(U_{10},U_{01}\) is trivial here,
so their frame terms vanish.

For the eight supports preceding \(Y\), the ones
that can meet these vertices are rows \(1,4,5,7,8\)
in the order of the displayed list.
Fix the global cocycle indicated for each surviving row.
\[
 \begin{array}{c|ccccc}
 \text{row}&1&4&5&7&8\\\hline
 \text{entry}&(1,2)&(0,3)&(0,1)&(1,2)&(0,3)
 \end{array}
\]
The image dimension of this coordinate map costs at most \(d+C\).
In the resulting kernel the two upper long-root
entries \((0,3),(1,2)\) vanish, either by this choice
or by the support.

In each row either \((0,1)\)
or \((1,0)\) also vanishes globally. The local
relation and \(b_v\ne0\) force both to vanish
locally. Thus every odd root entry is locally
a coboundary. The remaining global entries
have dimension at most \(10C\) by the strict
and even cohomology bounds. At most \(6g\)
frame coordinates remain, where \(g\le C\)
is the number of places over \(\ell\). Including
the three fixed cocycles, the tangent bound is
\[
                   2C+(d+C)+10C+6g\le d+20C .
\]
The completed support quotient therefore has
dimension at most \(d+20C+1\), with the extra
parameter from the coefficient DVR.

\proofstep{The support \(Y\).}
A representation supported in \(Y\) preserves
the two planes with indices \(\{0,3\}\) and
\(\{1,2\}\). Its two determinants are both the
similitude \(\nu\).
At any local point where the first lower flag
entry \((N_v)_{10}\) is nonzero, the ordinary first
vector has nonzero projection on each plane.

Let \(u_1,u_2\) be the two nonzero projections of
the first flag vector. Both satisfy
\(r(h)u_i=\psi_{v,0}(h)u_i\) on the decomposition group,
since the two plane projections commute with \(r(h)\).
The determinant \(\nu\) of each plane determines its other
graded character. Thus the full composition multiset is
\[
       \{\psi_{v,0},\psi_{v,0},\nu/\psi_{v,0},\nu/\psi_{v,0}\}.
\]
Since the pairing gives \(\psi_{v,3}=\nu/\psi_{v,0}\),
the multiset forces on the same decomposition group the coincidence
\[
               \psi_{v,1}=\psi_{v,0}\quad\text{or}\quad
               \psi_{v,1}=\psi_{v,3}.
\] At every prime of the completion at a
vertex in \(\mathcal A\) the frame entry stays
nonzero, since it is a unit at the vertex.

The dimension comparison with the characteristic
polynomial and ordinary character ring also holds
on the part of \(Y\) outside the first eight
supports. On this part each plane has both of
its off-diagonal entries nonzero somewhere.
Over the fraction field, the four eigenvalues
of the normalizer are algebraic over the base.
If \(\lambda_i\) are the marked eigenvalues, the four traces are
\[
    \operatorname{tr}(r(g_0)^j r(g))
            =\sum_i\lambda_i^j r(g)_{ii}\qquad(0\le j<4).
\]
Invert the Vandermonde matrix to read the diagonal entries. The same holds for
the diagonal entries of \(r(gh)\).
Subtracting products of the known diagonals
therefore determines
\[
                  r(g)_{ij}r(h)_{ji}
\]
for each opposite pair of allowed off-diagonal
positions and all \(g,h\).

Choose one nonzero
value in one direction for each plane. All other
entries in both directions in that plane are
then obtained from these products by division.
Thus two field generators, together with the
four entries determining each local lower
symplectic frame, suffice over an algebraic
extension of the base field. The finite type
dimension comparison increases the dimension
by at most \(2+4g\), just as at the absolutely
irreducible primes.

\proofstep{The selected arc.}

An arc avoiding the excluded loci lets us apply the one-point support result.

At an intersection of adjacent components through
an extra vertex, start on the modular component
and cut by the coefficient uniformizer. Choose
a prime in the intersection with quotient
dimension at least \(N=2d+1-2g\).
The depth inequalities give
\[
     d+20C+1<N,\qquad 2d+1-d_v+2+4g<N .
\]
If \(Z\) is the chosen integral closed subset, its
dimension is at least \(N\). The first inequality gives
\(Z\not\subseteq\) any of the eight small reducible supports.
The second gives the same exclusion for inertia coincidence loci,
using the dimension comparison on \(Y\).
In particular the prime cannot be supported
in \(Y\): its unit lower frame entry would
force a coincidence on the full decomposition
group and hence on the open inertia.
The prime therefore has absolutely irreducible
generic representation and separated inertia
characters.

In the local coordinate domain of \(Z\), choose \(0\ne f_i\) from each forbidden-locus ideal and \(0\ne t\) from the maximal ideal. The element
\[
 f=t\prod_i f_i\ne0
\]
belongs to the maximal ideal and vanishes on each forbidden locus. Lemma~\ref{lem:local-curve-avoidance} gives a one-dimensional quotient retaining that product. Normalizing as in the component argument gives the centered formal arc with generic point outside the exclusions. Local vanishing at this
arc puts the shared Hecke ideal in the adjacent
component. This proves the required implication
from one modular component to all modular
components at both new kinds of vertices.
It applies to every vertex in the residual
chain, so the backwards argument of
Paragraph~\ref{par:count-two-propagation} applies
once the integral seed lies in the ordinary
support with its cuspidal attachment.
\label{cor:ordinary-count-iii-components-1}
\end{proof}
